% Part:sets-functions-relations
% Chapter: sets
% Section: reduction

\documentclass[../../../include/open-logic-section]{subfiles}

\begin{document}

\olfileid{sfr}{siz}{red}

\olsection{Reduksi}

\begin{editorial}
  Bagean iki mbuktekake sipat ora bisa didhaptar kanthi reduksi,
  cocog karo asil ing \olref[nen]{sec}. Versi alternatif kang rada luwih
  ringkes, cocog karo asil ing \olref[nen-alt]{sec}, diwenehake ing
  \olref[red-alt]{sec}.
\end{editorial}

Kita wis nuduhake yen $\Pow{\PosInt}$ iku !!{nonenumerable} nganggo
argumen diagonalisasi. Sadurunge, kita wis duwe bukti yen $\Bin^\omega$,
himpunan kabeh rerangken tanpa wates saka $0$ lan $1$, iku !!{nonenumerable}.
Ana cara liya kanggo mbuktekake yen $\Pow{\PosInt}$ iku !!{nonenumerable}:
tuduhna yen \emph{yen $\Pow{\PosInt}$ iku !!{enumerable}, mula $\Bin^\omega$
  uga !!{enumerable}}. Amarga kita ngerti yen $\Bin^\omega$ ora
!!{enumerable}, $\Pow{\PosInt}$ uga ora bisa mangkono. Iki diarani
\emph{ngreduksi} sawijining masalah menyang masalah liyane. Ing kene,
masalah ngenumerasi $\Bin^\omega$ kita reduksi menyang masalah ngenumerasi
$\Pow{\PosInt}$. Solusi kanggo masalah kang pungkasan, yaiku enumerasi
$\Pow{\PosInt}$, bakal ngasilake solusi kanggo masalah kapisan, yaiku
enumerasi $\Bin^\omega$.

Kepriye carane ngreduksi masalah ngenumerasi himpunan~$B$ menyang
masalah ngenumerasi himpunan~$A$? Kita menehi cara kanggo ngowahi enumerasi
saka~$A$ dadi enumerasi saka~$B$. Cara kang paling gampang yaiku netepake
fungsi !!a{surjective} $f\colon A \to B$. Yen $x_1$, $x_2$,
\dots{} ngenumerasi~$A$, mula $f(x_1)$, $f(x_2)$, \dots{}
ngenumerasi~$B$. Ing kasus iki, kita nggoleki fungsi surjektif
$f\colon \Pow{\PosInt} \to \Bin^\omega$.

\begin{prob}
Tuduhna yen ana fungsi !!{injective} $g\colon B \to A$, lan
$B$~iku !!{nonenumerable}, mula~$A$ uga mangkono. Tuduhna carane
nggunakake~$g$ kanggo ngowahi enumerasi saka~$A$ dadi enumerasi saka~$B$.
\end{prob}

\begin{proof}[Bukti {\olref[nen]{thm:nonenum-pownat}} kanthi reduksi]
Upamakna $\Pow{\PosInt}$ iku !!{enumerable}, mula ana
enumerasi $Z_{1}$, $Z_{2}$, $Z_{3}$, \dots

Tetepna fungsi $f \colon \Pow{\PosInt} \to \Bin^\omega$ kanthi netepake
$f(Z)$ dadi rerangken $s$ saengga $s(n) = 1$ yen lan mung yen $n \in Z$,
lan $s(n) = 0$ yen ora mangkono. Jeneng metuane diseragamake amarga
indeks k ing sumber Inggris ora kaiket (OLSIZ-004). Iki cetha netepake sawijining fungsi,
amarga yen $Z \subseteq \PosInt$, saben $n \in \PosInt$ mesthi
!!a{element} saka $Z$ utawa dudu anggotane. Upamane, himpunan $2\PosInt = \{2,
4, 6, \dots\}$ saka wilangan genep positif dipetakake menyang rerangken
$010101\dots$, himpunan kosong dipetakake menyang $0000\dots$, lan
himpunan $\PosInt$ dhewe menyang $1111\dots$.

Fungsi iki uga !!{surjective}: saben rerangken $0$ lan $1$ cocog
karo sawijining himpunan wilangan wutuh positif, yaiku himpunan kang
anggotane wilangan kang nuduhake posisi rerangken kang isine~$1$.
Kanthi luwih pas, upamakna $s \in \Bin^\omega$. Tetepna $Z
\subseteq \PosInt$ kanthi:
\[
Z = \Setabs{n \in \PosInt}{s(n) = 1}
\]
Dadi $f(Z) = s$, kaya kang bisa dipriksa saka definisi~$f$.

Saiki gatekna dhaptar
\[
f(Z_1), f(Z_2), f(Z_3), \dots
\]
Amarga $f$ iku !!{surjective}, saben anggota $\Bin^\omega$ mesthi
katon minangka nilai~$f$ kanggo sawijining argumen, mula mesthi
katon ing dhaptar. Dadi dhaptar iki ngenumerasi kabeh~$\Bin^\omega$.

Mula yen $\Pow{\PosInt}$ iku !!{enumerable}, $\Bin^\omega$ mesthi
!!{enumerable}. Nanging $\Bin^\omega$ iku !!{nonenumerable}
(\olref[nen]{thm:nonenum-bin-omega}). Dadi $\Pow{\PosInt}$ iku
!!{nonenumerable}.
\end{proof}

\begin{explain}
Arah reduksi gampang mbingungake.
Upamane, fungsi !!a{surjective} $g \colon \Bin^\omega \to B$
\emph{ora} mbuktekake yen $B$ iku !!{nonenumerable}. (Gatekna $g
\colon \Bin^\omega \to \Bin$ kang ditetepake dening $g(s) = s(1)$,
yaiku fungsi kang metakake rerangken $0$ lan $1$ menyang !!{element} kapisane.
Fungsi iki !!{surjective}, amarga ana rerangken kang diwiwiti $0$ lan
ana kang diwiwiti $1$. Nanging $\Bin$ winates.) Gatekna uga yen fungsi~$f$
kudu !!{surjective}; yen ora, argumen iki ora lumaku:
$f(x_1)$, $f(x_2)$, \dots{} ora mesthi ngemot
kabeh !!{element}s saka~$B$. Upamane,
\[
h(n) = \underbrace{00\dots0}_{\text{$n$ wilangan $0$}}111\dots
\]
netepake fungsi $h\colon \PosInt \to
\Bin^\omega$, nanging $\PosInt$ iku !!{enumerable}. Buntut siji tanpa
wates ndadekake saben nilai kalebu kodomain; sumber Inggris mung
nampilake string nol winates (OLSIZ-005).
\end{explain}

\begin{prob}
Tuduhna yen himpunan kabeh \emph{himpunan saka} pasangan wilangan wutuh positif
iku !!{nonenumerable} nganggo argumen reduksi.
\end{prob}

\begin{prob}\label{sfr:siz:red:prob:nat-nat}
  Tuduhna yen himpunan~$X$ saka kabeh fungsi $f\colon \Nat \to \Nat$
  iku !!{nonenumerable} nganggo argumen reduksi. (Pituduh: wenehana fungsi
  surjektif saka $X$ menyang~$\Bin^\omega$.)
\end{prob}

\begin{prob}
Tuduhna yen $\Nat^\omega$, himpunan rerangken tanpa wates saka
wilangan natural, iku !!{nonenumerable} nganggo argumen reduksi.
\end{prob}

\begin{prob}
Upamakna $P$ himpunan fungsi saka himpunan wilangan wutuh positif
menyang himpunan $\{0\}$, lan $Q$ himpunan fungsi \emph{parsial}
saka himpunan wilangan wutuh positif menyang himpunan $\{0\}$.
Tuduhna yen $P$~iku !!{enumerable} lan $Q$~ora. (Pituduh: reduksinen
masalah ngenumerasi $\Bin^\omega$ menyang ngenumerasi~$Q$.)
\end{prob}

\begin{prob}
Upamakna $S$ himpunan kabeh fungsi !!{surjective} saka himpunan
wilangan wutuh positif menyang himpunan \{0,1\}, yaiku $S$ ngemot kabeh
!!{surjective}~$f\colon \PosInt \to \Bin$. Tuduhna yen $S$ iku
!!{nonenumerable}.
\end{prob}

\begin{prob}
Tuduhna yen himpunan~$\Real$ saka kabeh wilangan real iku !!{nonenumerable}.
\end{prob}

\end{document}
